% !TEX TS-program = lualatex
Выбор языка программирования высокого уровня для реализации поставленной задачи был обусловлен спецификой выбранного алгоритмического решения. В качестве основного инструмента был выбран Python --- мультипарадигменный язык программирования высокого уровня, который позволяет лаконично и выразительно описывать сложные программные и математические конструкции. Благодаря широкому набору библиотек, активному сообществу и простоте синтаксиса, Python занимает лидирующие позиции в области искусственного интеллекта, машинного обучения и анализа данных на сегодняшний день \cite{programming-languages-rating}.

Проект, представленный в данной работе, базируется на версии Python 3.13. Управление средой разработки и зависимостями осуществляется с помощью менеджера версий uv, обеспечивающего быстрый и изолированный запуск приложений, а также удобное управление пакетами.

Для обеспечения высокого качества и стабильности кода применяются инструменты статического анализа. В частности, используется Pyright --- быстрый и точный синтаксический анализатор, ориентированный на проверку типизации кода на Python и доступный для редактора кода VSCode. Дополнительно применяется Ruff --- высокопроизводительный линтер, отвечающий за соответствие кода современным стилевым стандартам, обнаружение потенциальных ошибок и автоматическое исправление распространённых нарушений.

Для построения, обучения и валидации интеллектуальных агентов на основе нейронных сетей используется библиотека TensorFlow, разработанная Google. TensorFlow предоставляет мощный и гибкий инструментарий для реализации нейронных сетей различной сложности. Одним из ключевых преимуществ TensorFlow является его модульность и поддержка как низкоуровневого, так и высокоуровневого API \cite{tensorflow-api}, что позволяет детально управлять вычислительным графом при необходимости.

В качестве высокоуровневого интерфейса к TensorFlow применяется Keras --- встроенный модуль, обеспечивающий интуитивно понятный способ описания архитектур нейронных сетей \cite{keras-api}. Keras позволяет легко создавать и конфигурировать последовательные и функциональные модели, а также упрощает процессы компиляции, обучения и оценки моделей благодаря лаконичному синтаксису и обширной документации.

Для организации взаимодействия между сервером игры и моделью используется современный веб-фреймворк FastAPI, который предоставляет асинхронную, высокопроизводительную архитектуру на базе стандарта ASGI (Asynchronous Server Gateway Interface), что делает его оптимальным выбором для создания RESTful API в проектах с высокой нагрузкой и требующих низкой задержки отклика. Благодаря встроенной поддержке валидации данных с использованием Pydantic и автоматической генерации документации OpenAPI, FastAPI существенно упрощает разработку, тестирование и сопровождение веб-сервисов \cite{fastapi-api}.

Для валидации и сериализации данных в проекте используется библиотека Pydantic, которая обеспечивает строгую типизацию и удобное управление структурами данных. Pydantic основана на использовании аннотаций типов Python и позволяет автоматически проверять корректность входных данных, преобразовывать их и создавать валидные объекты моделей \cite{pydantic-api}.

Криптографическая библиотека PyCryptodome применяется для обеспечения конфиденциальности и безопасности данных. Она представляет собой низкоуровневую реализацию современных криптографических примитивов, поддерживающих симметричное и асимметричное шифрование, генерацию случайных чисел, цифровые подписи и хэширование \cite{pycryptodome-api}. В рамках проекта PyCryptodome используется для безопасного хранения и передачи конфиденциальной информации об игровой сессии, а также для реализации механизмов контроля целостности и аутентификации сообщений.

Для реализации интерфейса игры клиентская часть приложения использует технологии HTML5, CSS и JavaScript. HTML5 выступает в роли структурного каркаса веб-страниц, предоставляя семантически богатый набор тегов и элементов для корректного отображения контента и мультимедийных данных.

Cascading Style Sheets (CSS) применяются для стилизации визуальных компонентов интерфейса, позволяя задавать оформление и обеспечивать оптимальный внешний вид приложения на различных устройствах и экранах. В проекте используется современный подход к построению интерфейсов с применением Flexbox и Grid Layout, что способствует гибкости и масштабируемости дизайна.

JavaScript играет ключевую роль в реализации динамической логики на стороне клиента, обеспечивая обработку пользовательских событий, асинхронное взаимодействие с сервером и обновление интерфейса без полной перезагрузки страницы. Для упрощения разработки и повышения производительности применяется современный стандарт ECMAScript 6 и выше, который включает модули, стрелочные функции, промисы и async/await.
